\documentclass[11pt]{article}

\usepackage[a4paper,margin=1in]{geometry}
\usepackage{amsmath,amssymb}
\usepackage{hyperref}

\setlength{\parindent}{0pt}
\setlength{\parskip}{0.7em}

\title{Instruction: Implement Generalized Three-Level Squeezing Parameter}
\author{}
\date{}

\begin{document}
\maketitle

Construct the generalized three-level squeezing parameter for the custom MCWF
code. The goal is to compute \(\xi^2_{\rm gen}(t)\) at each saved timestep after
the simulation has run.

The squeezing calculation should preferably be implemented as a standalone
post-processing function that takes the simulation result or observable output as
input. If it is simpler, it may be integrated into the observable function, but
it must be optional via a boolean flag, since this calculation may be expensive
and is not always needed.

If additional per-timestep data must be saved for this calculation, add a
boolean flag such as \texttt{save\_squeezing\_data} or
\texttt{compute\_generalized\_squeezing}.

\section*{Per-timestep calculation}

At each timestep, do the following.

\subsection*{1. Define the effective dressed state \(|1\rangle\)}

Work in the single-particle basis

\begin{equation}
(|u\rangle, |d\rangle, |e\rangle).
\end{equation}

Use the ansatz

\begin{equation}
|1\rangle =
0|u\rangle
+
\cos\frac{\theta_J}{2}|d\rangle
+
e^{-i\phi_J}\sin\frac{\theta_J}{2}|e\rangle.
\end{equation}

Here \(\theta_J,\phi_J\) should be found by comparing this ansatz to the
simulation state at the current timestep. Follow the active-manifold
Bloch-vector averaging convention in:

\begin{verbatim}
docs/instructions/bloch_vector_averaging.tex
\end{verbatim}

Reuse \texttt{active\_manifold\_angles} where appropriate.

\subsection*{2. Define the instantaneous mean direction \(|c\rangle\)}

Use the ansatz

\begin{equation}
|c\rangle
=
\cos\frac{\theta_S}{2}|u\rangle
+
e^{-i\phi_S}\sin\frac{\theta_S}{2}|1\rangle.
\end{equation}

Here \(\theta_S,\phi_S\) should be found by comparing this ansatz to the
instantaneous effective \(S\)-Bloch vector of the simulation state. The angles
\(\theta_J,\phi_J\) appearing inside \(|1\rangle\) are the angles found in step
1.

In the explicit \((u,d,e)\) basis,

\begin{equation}
|c\rangle =
\begin{pmatrix}
\cos(\theta_S/2) \\
e^{-i\phi_S}\sin(\theta_S/2)\cos(\theta_J/2) \\
e^{-i(\phi_S+\phi_J)}\sin(\theta_S/2)\sin(\theta_J/2)
\end{pmatrix}.
\end{equation}

\subsection*{3. Define the \(J\)-fluctuation direction \(|j\rangle\)}

Using the \(J\)-angles found in step 1, construct the state orthogonal to
\(|1\rangle\) inside the \((d,e)\) manifold:

\begin{equation}
|j\rangle =
0|u\rangle
-
\sin\frac{\theta_J}{2}|d\rangle
+
e^{-i\phi_J}\cos\frac{\theta_J}{2}|e\rangle.
\end{equation}

In vector form:

\begin{equation}
|j\rangle =
\begin{pmatrix}
0 \\
-\sin(\theta_J/2) \\
e^{-i\phi_J}\cos(\theta_J/2)
\end{pmatrix}.
\end{equation}

\subsection*{4. Define the \(S\)-fluctuation direction \(|s\rangle\)}

Using the \(S\)-angles from step 2 and the \(J\)-angles from step 1, construct

\begin{equation}
|s\rangle =
-\sin\frac{\theta_S}{2}|u\rangle
+
e^{-i\phi_S}\cos\frac{\theta_S}{2}|1\rangle.
\end{equation}

In the \((u,d,e)\) basis,

\begin{equation}
|s\rangle =
\begin{pmatrix}
-\sin(\theta_S/2) \\
e^{-i\phi_S}\cos(\theta_S/2)\cos(\theta_J/2) \\
e^{-i(\phi_S+\phi_J)}\cos(\theta_S/2)\sin(\theta_J/2)
\end{pmatrix}.
\end{equation}

The three states \(|c\rangle,|j\rangle,|s\rangle\) should form an orthonormal
single-particle basis up to numerical precision.

\subsection*{5. Construct the four local fluctuation operators}

Construct the four \(3\times3\) single-particle operators:

\begin{equation}
o_1 =
\frac{|c\rangle\langle j|+|j\rangle\langle c|}{2},
\end{equation}

\begin{equation}
o_2 =
\frac{|c\rangle\langle j|-|j\rangle\langle c|}{2i},
\end{equation}

\begin{equation}
o_3 =
\frac{|c\rangle\langle s|+|s\rangle\langle c|}{2},
\end{equation}

\begin{equation}
o_4 =
\frac{|c\rangle\langle s|-|s\rangle\langle c|}{2i}.
\end{equation}

Then construct the corresponding collective operators

\begin{equation}
O_a = \sum_i o_a^{(i)},
\qquad a=1,2,3,4.
\end{equation}

Do not construct full \(3^N\)-dimensional tensor-product operators unless
absolutely necessary. Use the existing reduced \((N_J,n_e)\) basis and exploit
symmetry. Each collective operator should be represented as a sparse operator
acting on the reduced simulation basis.

If possible, construct \(O_a\) from precomputed collective one-body transition
operators

\begin{equation}
A_{\mu\nu}=\sum_i |\mu_i\rangle\langle \nu_i|,
\qquad
\mu,\nu\in\{u,d,e\}.
\end{equation}

Then

\begin{equation}
O_a =
\sum_{\mu,\nu}
(o_a)_{\mu\nu} A_{\mu\nu}.
\end{equation}

This avoids building large tensor-product matrices.

\subsection*{6. Construct the covariance matrix}

For the current state \(|\psi(t)\rangle\), compute

\begin{equation}
\mu_a = \langle O_a\rangle.
\end{equation}

Then construct the \(4\times4\) covariance matrix

\begin{equation}
C_{ab}
=
\frac{1}{2}
\langle O_aO_b+O_bO_a\rangle
-
\mu_a\mu_b.
\end{equation}

For efficiency, avoid explicitly constructing \(O_aO_b\). Instead use

\begin{equation}
C_{ab}
=
\mathrm{Re}\left[
\langle O_a\psi | O_b\psi\rangle
\right]
-
\mu_a\mu_b.
\end{equation}

\subsection*{7. Minimum fluctuation direction}

Diagonalize the \(4\times4\) covariance matrix and take

\begin{equation}
\lambda_{\min}(C).
\end{equation}

This is the minimum generalized transverse variance.

\subsection*{8. Generalized squeezing parameter}

Compute

\begin{equation}
\xi^2_{\rm gen}(t)
=
\frac{N\lambda_{\min}(C)}
{\langle N_c/2\rangle^2},
\end{equation}

where

\begin{equation}
N_c =
\sum_i |c_i\rangle\langle c_i|.
\end{equation}

Compute \(\langle N_c\rangle\) using the same collective-operator construction:

\begin{equation}
N_c =
\sum_{\mu,\nu}
(|c\rangle\langle c|)_{\mu\nu} A_{\mu\nu}.
\end{equation}

If the state is well polarized along \(|c\rangle\), then
\(\langle N_c/2\rangle\approx N/2\), but the code should compute it explicitly
rather than assuming this.

\section*{Data that may need to be saved per timestep}

Save enough data to reconstruct the squeezing calculation after the simulation:

\begin{itemize}
\item the full state vector \(|\psi(t)\rangle\) in the reduced \((N_J,n_e)\)
  basis, or equivalent data sufficient to reconstruct it;
\item the timestep values \(t\);
\item \(N\), \(\Gamma\), \(\Omega\), \(\delta\), and any sector truncation
  metadata;
\item the basis/index mapping for \((N_J,n_e)\);
\item if not recomputed during post-processing, save
  \(\theta_J(t),\phi_J(t),\theta_S(t),\phi_S(t)\).
\end{itemize}

Prefer recomputing angles from the state during post-processing if this avoids
storing redundant data.

\end{document}
